\section{Sim2Real, Real2Sim y datos sintéticos}

Esta sección explica primero los términos más frecuentes, porque los juicios posteriores sobre la industria se apoyan en estos conceptos. Sim2Real es ir de la simulación a la realidad; Real2Sim es trasladar la realidad a la simulación; Synthetic Data son datos aptos para entrenamiento generados mediante simulación, algoritmos y procesos programáticos. Los tres no son buzzwords paralelas, sino posiciones distintas dentro de un mismo ciclo de ingeniería: el mundo real plantea el problema, la simulación lo amplía, el modelo se entrena con datos sintéticos y los resultados del despliegue vuelven a calibrar la simulación.

La experiencia de Xie Chen (谢晨) en Cruise aporta una lección decisiva. Al principio, el equipo recurrió a ingenieros de videojuegos para recrear San Francisco; el resultado parecía Matrix, pero el área de percepción no podía usar esos datos directamente, porque no se habían definido según objetivos de aprendizaje automático. Lo que después funcionó de verdad fue construir primero un sistema de evaluación que determinara si los datos sintéticos mejoraban el modelo y, a partir de ello, modificar los recursos, los escenarios, el ruido de los sensores y la cadena de generación de datos.

\begin{knowledgebox}{Términos clave: simulación, datos sintéticos y transferencia}
\begin{tabularx}{\textwidth}{p{0.20\textwidth}p{0.35\textwidth}X}
\toprule
Término & Significado básico & Cuestión central en este episodio \\
\midrule
Sim2Real & Entrenar o generar datos en simulation y luego transferirlos al mundo real & El éxito del despliegue real determina si la simulación sirve. \\
Sim2Real gap & Diferencia de distribución entre la simulación y la realidad & No se busca una diferencia nula, sino determinar a partir de qué magnitud de diferencia los datos empiezan a ser útiles. \\
Synthetic Data & Datos generados mediante simulación, algoritmos, procesos programáticos o colaboración humano-máquina & Lo decisivo no es la cantidad generada, sino la densidad de información, la calidad y la utilidad para el modelo. \\
Real2Sim & Trasladar a la simulación escenarios reales, recursos, parámetros físicos y el propio cuerpo del robot & La reconstrucción visual es relativamente sencilla; lo difícil son los parámetros físicos y la mecánica de la interacción. \\
Evaluation Loop & Inferir la calidad de los datos a partir del entrenamiento del modelo y de los resultados del despliegue real & Sin un sistema de evaluación, al equipo de datos le resulta difícil saber qué corregir. \\
\bottomrule
\end{tabularx}
\end{knowledgebox}

\subsection{De un juguete vistoso a datos para aprendizaje automático}

Esta subsección responde una pregunta concreta: por qué una simulación que «parece real» no equivale a una simulación «útil para el modelo». El sistema de simulación inicial de Cruise podía recrear manzanas enteras de la ciudad, pero si el ruido de los sensores, la distribución de la iluminación, las anotaciones, la probabilidad de los escenarios y los eventos de cola larga no correspondían a lo que requería el entrenamiento del modelo, el sistema de percepción no obtenía ningún beneficio. El primer paso de Xie Chen no fue mejorar la calidad gráfica, sino construir un sistema de evaluación.

El sistema de evaluación tiene dos tipos de indicadores. El primero es el realismo absoluto: si la luz, el color, el ruido de los sensores, la geometría, las anotaciones y la distribución de escenarios coinciden con los datos reales. El segundo es la utilidad: si, tras alimentar al modelo con este conjunto de datos sintéticos, su desempeño en la tarea objetivo mejora. El ponente lo compara con un espresso frente a un café helado rebajado con agua: con el tiempo, los datos se valoran cada vez más por «cuánta información adicional aporta al modelo cada unidad de datos», y no solo por su volumen.

\begin{importantbox}{El sistema de evaluación va primero}
Lo primero que debe preguntarse un equipo de datos sintéticos no es «cómo generar más», sino «qué datos hacen que el modelo mejore». Sin un sistema de evaluación, el equipo de simulación solo puede optimizar la apariencia visual; con él, el equipo puede saber si debe corregir el ruido de los sensores, la distribución de escenarios, los parámetros físicos de los recursos o la cadena de control de calidad de los datos.
\end{importantbox}

La utilidad de los datos puede expresarse como un objetivo simplificado:
\[
U(D_s)=\Delta M_{\mathrm{target}}-\lambda C(D_s)-\gamma G(D_s,D_r)
\]
donde $D_s$ representa los datos sintéticos, $D_r$ los datos reales, $\Delta M_{\mathrm{target}}$ la mejora en el indicador del modelo objetivo, $C(D_s)$ el costo de generación y control de calidad, $G(D_s,D_r)$ el gap relevante entre lo sintético y lo real, y $\lambda$ y $\gamma$ son los pesos del costo y del gap. Esta fórmula no procede textualmente de la entrevista; es una formulación didáctica de su metodología: los datos sintéticos eficaces deben considerar a la vez el beneficio para el modelo, el costo de producción y la diferencia con la realidad.

\subsection{Flujo de datos sintéticos y proporciones}

A continuación, esta subsección aterriza el flujo en un ejemplo: un peatón que aparece de improviso en la calle. Este tipo de corner case es escaso en los datos reales, pero crucial para la seguridad. El flujo de datos sintéticos no consiste en copiar un video diez mil veces, sino en generar variantes sistemáticas en torno a un problema real: distintas edades, complexiones, ropa, trayectorias, carriles, clima, iluminación, condiciones de tráfico y ruido de los sensores. Después, el algoritmo se ejecuta diez mil veces en la simulación para descubrir si la percepción, la predicción y la planificación siguen fallando.

No existe una proporción fija. En el conjunto total de datos de la conducción autónoma, la flota real genera a diario una gran cantidad de datos reales, de modo que los datos sintéticos quizá no sean la mayor parte del total; pero para los corner cases escasos, las muestras reales pueden no bastar ni en un año, por lo que a un solo problema real se le asocian miles o decenas de miles de variantes sintéticas. La experiencia que se comparte en el programa es que, en conjunto, en algún momento cerca del 30\% de los datos fueron sintéticos, mientras que en los escenarios de cola larga la proporción real/sintético puede superar 1:99 o 1:100.

\begin{knowledgebox}{El flujo paso a paso: de un problema a un conjunto de datos de entrenamiento}
\begin{tabularx}{\textwidth}{p{0.20\textwidth}p{0.35\textwidth}X}
\toprule
Paso & Qué se hace & Por qué importa \\
\midrule
Localización del problema real & Encontrar un caso de peligro o de falla & La generación de datos debe partir de riesgos reales. \\
Parametrización del escenario & Variar peatones, carriles, clima, iluminación y flujo de tráfico & Permite que el modelo vea múltiples manifestaciones de un mismo problema. \\
Simulación de sensores & Simular cámaras, radares, ruido y perspectivas de observación & Lo que el modelo recibe al final son datos de sensores. \\
Anotación automática y control de calidad & Generar etiquetas 2D/3D/4D y verificar su calidad & Los errores de anotación contaminan directamente el entrenamiento. \\
Retroalimentación al modelo & Usar los datos para entrenar o evaluar percepción/predicción/planificación & La utilidad debe demostrarse con los indicadores del modelo y los resultados del despliegue. \\
\bottomrule
\end{tabularx}
\end{knowledgebox}

\begin{warningbox}{La proporción no puede discutirse al margen de la tarea}
Decir «1\% de datos reales y 99\% de datos sintéticos» se presta a confusión. La pregunta correcta es: ¿para qué tarea, qué modo de falla, qué módulo del modelo y qué etapa del despliegue tienen los datos sintéticos una utilidad marginal mayor? Las proporciones del conjunto total de datos, de los escenarios de cola larga y de los entornos de aprendizaje por refuerzo pueden ser completamente distintas.
\end{warningbox}

\subsection{Resumen de la sección}

Sim2Real y Real2Sim forman un ciclo cerrado de datos. El valor de los datos sintéticos no lo determina el realismo de la imagen, sino dos tipos de evaluación: si coinciden con las variables decisivas del mundo real y si hacen que el modelo objetivo mejore en tareas reales.
